% ============================================================================
% 04_scaling_law / 01_scaling_law
% Section body for inclusion immediately after the corresponding \section
% command in the master document. Do not compile this file independently.
% ============================================================================
%
% This file intentionally contains no document-level commands, package imports,
% \chapter, \section, or \printbibliography commands. The master document is
% responsible for those elements.
%
% Figure source fragments are stored in the root-level figures/ directory.
% References are stored in the root-level bibliography.bib file.
%
% ----------------------------------------------------------------------------
% Opening transition and conceptual motivation
% ----------------------------------------------------------------------------

The preceding discussions of language-model architecture, optimization, and
training objectives establish how an autoregressive model is parameterized and
how its parameters are adjusted to improve next-token prediction. A further
question follows naturally from these foundations: how should the size of the
model, the quantity of training data, and the computational resources be
increased as the intended system becomes larger? Answering this question
requires more than observing that larger models often achieve lower predictive
loss. It requires a quantitative account of how performance changes with scale,
how the different resources interact, and how a limited computational budget
should be allocated.

Training a large language model is expensive not only because the final model
may contain billions or trillions of parameters, but also because identifying
a suitable configuration requires experiments with different architectures,
training-data quantities, optimization settings, and training durations.
Exhaustively testing every configuration at the target scale is generally
impractical. Scaling laws address this difficulty by fitting mathematical
relationships to observations obtained from a manageable collection of
experiments and using those relationships to predict the behavior of larger
systems. They therefore connect empirical machine-learning practice with
statistical estimation, numerical optimization, and constrained resource
allocation \parencite{Sengupta2025ScalingSurvey,Li2025MisFitting}.

For autoregressive language models, the central object is usually a predictive
loss measured on a specified data distribution. The scaling variables
traditionally include the number of model parameters, the number of training
tokens, and the computational expenditure required to process those tokens.
The conventional framework is subsequently extended when inference latency,
architectural shape, or sparse conditional computation becomes important.
These extensions are necessary because a model that is optimal in pre-training
FLOPs need not be optimal in memory use, inference latency, throughput, or
deployment cost.

A scaling law is an empirical predictive model, not a universal identity.
Its coefficients and exponents depend on the model family, data distribution,
training procedure, evaluation metric, experimental range, and fitting
methodology. The objective of this section is to develop the foundations in a
logical order: first defining the predictive quantity and scale variables,
then deriving power-law and compute-optimal relationships, explaining how
these relationships are estimated and validated, and finally examining
inference-aware, architecture-aware, and mixture-of-experts (MoE) extensions.

% ----------------------------------------------------------------------------
% Main exposition
% ----------------------------------------------------------------------------

\subsection{Predictive loss and the fundamental scaling variables}
\label{subsec:scaling-foundations-loss}

\paragraph{Autoregressive prediction and cross-entropy.}
An autoregressive language model factorizes the probability of a token
sequence into conditional next-token probabilities. For a sequence
\[
    X=(x_1,x_2,\ldots,x_T),
\]
where \(T\) is the sequence length, the model with parameters
\(\boldsymbol{\theta}\) represents
\begin{equation}
    p_{\boldsymbol{\theta}}(X)
    =
    \prod_{t=1}^{T}
    p_{\boldsymbol{\theta}}(x_t\mid x_{<t}),
    \label{eq:scaling-autoregressive-factorization}
\end{equation}
where \(x_{<t}=(x_1,\ldots,x_{t-1})\) denotes the context preceding position
\(t\). This factorization follows from the chain rule of probability and
reduces sequence modeling to a succession of conditional prediction problems.

Let the vocabulary contain \(V\) tokens. At position \(t\), the model produces
a probability vector
\[
    \mathbf{p}^{(t)}_{\boldsymbol{\theta}}
    =
    \left[
    p_{\boldsymbol{\theta}}(1\mid x_{<t}),\ldots,
    p_{\boldsymbol{\theta}}(V\mid x_{<t})
    \right]^{\mathsf T}
    \in\mathbb{R}^{V}.
\]
The target token can be represented by a one-hot vector
\(\mathbf{y}^{(t)}\in\{0,1\}^{V}\), whose component corresponding to the
observed token is one. The cross-entropy at this position is
\begin{equation}
    \ell_t(\boldsymbol{\theta})
    =
    -\sum_{v=1}^{V}
    y_v^{(t)}
    \log p_{\boldsymbol{\theta}}(v\mid x_{<t})
    =
    -\log p_{\boldsymbol{\theta}}(x_t\mid x_{<t}).
    \label{eq:scaling-token-cross-entropy}
\end{equation}
The second equality follows from the one-hot structure of the target
distribution.

For \(D\) token-prediction targets, the empirical average negative
log-likelihood is
\begin{equation}
    \widehat{\mathcal{L}}_{\mathrm{train}}
    (\boldsymbol{\theta})
    =
    -\frac{1}{D}
    \sum_{t=1}^{D}
    \log p_{\boldsymbol{\theta}}(x_t\mid x_{<t}).
    \label{eq:scaling-empirical-loss}
\end{equation}
Minimizing this quantity encourages the model to assign higher probability
to the observed next tokens. At the population level, the corresponding
predictive risk is an expectation over the data-generating distribution
\(\mathcal{D}\):
\begin{equation}
    \mathcal{L}_{\mathcal{D}}(\boldsymbol{\theta})
    =
    \mathbb{E}_{X\sim\mathcal{D}}
    \left[
    -\frac{1}{T}
    \sum_{t=1}^{T}
    \log p_{\boldsymbol{\theta}}(x_t\mid x_{<t})
    \right].
    \label{eq:scaling-population-loss}
\end{equation}
The empirical training loss and the population-level risk are not
interchangeable. Training loss is computed on observations used for
optimization, whereas validation or test loss is computed on separate
observations to estimate generalization. A scaling study must therefore state
which loss is being fitted and must not silently treat training and evaluation
measurements as equivalent.

For natural-log cross-entropy, perplexity is
\begin{equation}
    \operatorname{PPL}
    =
    \exp(\mathcal{L}).
    \label{eq:scaling-perplexity}
\end{equation}
Because the exponential function is strictly increasing, minimizing
cross-entropy and minimizing perplexity induce the same ordering of models,
although their numerical scales differ. Cross-entropy is particularly useful
for scaling analysis because it is a continuous measure of probabilistic
prediction. A discrete downstream metric, such as accuracy, can remain
unchanged over a range of loss improvements and then change abruptly.

This distinction matters when extrapolating pre-training behavior to
downstream tasks. A reduction in language-model loss does not imply a fixed or
universal increase in task accuracy. The relationship between loss and
accuracy can vary by task, as illustrated in the downstream comparisons
reported by \textcite{Bian2025InferenceEfficient}. Accordingly, a scaling law
must always be interpreted as a model of the metric that was actually fitted.

\paragraph{Model size.}
Let
\begin{equation}
    N=|\boldsymbol{\theta}|
    \label{eq:scaling-parameter-count}
\end{equation}
denote the number of trainable scalar parameters under a specified counting
convention. Parameters are the learned numerical entries of weight matrices,
bias vectors, embedding matrices, and other trainable tensors. In a
decoder-only Transformer, the parameter count depends on such architectural
quantities as hidden dimension, number of layers, attention projections, and
feed-forward dimensions.

The phrase ``number of parameters'' must be defined carefully. Some studies
count embedding and output-projection parameters, whereas others report
non-embedding parameters. A useful bookkeeping identity is
\begin{equation}
    N_{\mathrm{total}}
    =
    N_{\mathrm{nonembed}}+N_{\mathrm{embed}},
    \label{eq:scaling-parameter-count-convention}
\end{equation}
where the exact components included in \(N_{\mathrm{embed}}\) depend on the
study's convention. These alternatives can assign different numerical sizes
to the same architecture and can therefore change fitted scaling coefficients
and resource ratios. The distinction becomes even more important for sparse
models, in which stored parameters and parameters activated for a token are
not equal \parencite{Li2025MisFitting,Ludziejewski2025JointMoE}.

\paragraph{Training-data volume.}
Let \(D\) denote the number of training tokens processed during pre-training.
The token count depends on the tokenizer and preprocessing pipeline; a fixed
quantity of raw text can yield different token counts under different
tokenization procedures. It is therefore necessary to distinguish raw corpus
size from tokenized training volume.

It is also useful to distinguish unique tokens from processed token
exposures. If a corpus contains \(D_{\mathrm{unique}}\) tokens and is traversed
for \(E_{\mathrm{epoch}}\) complete epochs, then, under the simplifying
assumption that every epoch processes the entire corpus,
\begin{equation}
    D_{\mathrm{processed}}
    =
    E_{\mathrm{epoch}}D_{\mathrm{unique}}.
    \label{eq:scaling-repeated-data}
\end{equation}
The processed-token count measures training exposures; it does not measure
the quantity of unique information available to the model. Repeated exposure
to the same observations can behave differently from adding new data.
Moreover, datasets with equal token counts need not have equal statistical
value: domain composition, repetition, noise, and information content can
affect the learning outcome. Data volume is therefore a useful scale variable,
but not a complete description of data quality
\parencite{Sengupta2025ScalingSurvey,Li2025MisFitting}.

\paragraph{Training computation.}
Let \(C\) denote training computation, commonly approximated in floating-point
operations (FLOPs). A FLOP is a floating-point arithmetic operation. FLOPs
approximate arithmetic work but do not directly determine wall-clock time,
energy use, monetary cost, or hardware utilization.

For a conventional dense Transformer, a widely used approximation is
\begin{equation}
    C(N,D)\approx 6ND.
    \label{eq:scaling-six-nd}
\end{equation}
The factor of approximately six is a conventional accounting of forward and
backward arithmetic per parameter per processed token: roughly two operations
per parameter for the forward computation and approximately four additional
operations for the backward computation. Thus, the approximate training
computation per token is \(6N\), and processing \(D\) tokens gives
\(C\approx 6ND\).

For example, a dense model with \(N=10^9\) parameters trained on
\(D=20\times10^9\) tokens has approximate training computation
\begin{align}
    C
    &\approx 6(10^9)(20\times10^9) \nonumber\\
    &=1.2\times10^{20}\ \text{FLOPs}.
    \label{eq:scaling-flops-example}
\end{align}
This estimate is an approximation, not an exact accounting of every
implementation. Sequence-length-dependent attention computation, embedding
operations, architectural details, and implementation-specific work can
introduce additional terms. Long-context models are a particularly important
case because attention-related computation can depend explicitly on context
length \parencite{Li2025MisFitting}.

The three conventional variables are therefore linked by
\begin{equation}
    C\approx 6ND.
    \label{eq:scaling-resource-identity}
\end{equation}
Under a fixed compute budget, increasing \(N\) requires decreasing \(D\);
increasing \(D\) requires either decreasing \(N\) or increasing \(C\). The
central resource-allocation question is consequently not whether parameters
or data are independently beneficial, but how they should be balanced to
minimize predictive loss under a specified computational constraint.

\begin{figure}[htbp]
    \centering
    % Built-in LaTeX picture-environment figure fragment; no TikZ dependency.
\begingroup
\setlength{\unitlength}{0.8mm}
\begin{picture}(150,58)
    \put(4,39){\framebox(34,13){\shortstack{Model size\\\(N\)}}}
    \put(4,10){\framebox(34,13){\shortstack{Processed tokens\\\(D\)}}}
    \put(58,25){\framebox(42,17){\shortstack{Training compute\\\(C\approx 6ND\)}}}
    \put(112,25){\framebox(34,17){\shortstack{Predictive loss\\\(L(N,D)\)}}}
    \put(38,45){\vector(1,0){20}}
    \put(38,16){\vector(1,1){20}}
    \put(100,33){\vector(1,0){12}}
    \put(5,2){\makebox(94,5)[l]{Dense-model approximation; architecture and data quality also matter.}}
\end{picture}
\endgroup

    \caption{The conventional scaling variables and their relationship.
    Model size \(N\) and processed training tokens \(D\) determine approximate
    dense-model training computation \(C\approx 6ND\), while predictive loss
    is modeled as a function of these variables.}
    \label{fig:scaling-resource-map}
\end{figure}

\subsection{Power-law formulations and compute-optimal allocation}
\label{subsec:scaling-power-laws}

\paragraph{Univariate power laws.}
A power law describes a relationship of the form
\begin{equation}
    y=ax^k,
    \label{eq:scaling-basic-power-law}
\end{equation}
where \(a\) is a positive coefficient and \(k\) is a scaling exponent. When
\(y\) is a loss that decreases with scale, a common form is
\begin{equation}
    L(x)=L_{\infty}+ax^{-\alpha},
    \qquad a>0,\quad \alpha>0,
    \label{eq:scaling-univariate-loss}
\end{equation}
where \(L_{\infty}\) is the asymptotic component under the assumed model and
\(\alpha\) controls the rate at which the scale-dependent excess loss
decreases.

Differentiating with respect to \(x\) gives
\begin{equation}
    \frac{dL}{dx}
    =
    -a\alpha x^{-\alpha-1}<0,
    \label{eq:scaling-power-first-derivative}
\end{equation}
so increasing \(x\) decreases the modeled loss. The second derivative is
\begin{equation}
    \frac{d^2L}{dx^2}
    =
    a\alpha(\alpha+1)x^{-\alpha-2}>0,
    \label{eq:scaling-power-second-derivative}
\end{equation}
which expresses diminishing absolute improvements as scale increases.

Subtracting the asymptotic term and taking logarithms yields
\begin{equation}
    \log(L-L_{\infty})
    =
    \log a-\alpha\log x.
    \label{eq:scaling-log-linear}
\end{equation}
The slope in log-log coordinates is \(-\alpha\). Equivalently,
\[
    \frac{d\log(L-L_{\infty})}{d\log x}=-\alpha,
\]
so the exponent is the elasticity of excess loss with respect to scale. If
the scale is multiplied by \(q>0\), then
\begin{equation}
    \frac{L(qx)-L_{\infty}}{L(x)-L_{\infty}}
    =
    q^{-\alpha}.
    \label{eq:scaling-power-ratio}
\end{equation}
Thus, the exponent quantifies the proportional reduction in excess loss
associated with a proportional increase in the scaling variable.

\paragraph{Multivariable power laws.}
Language-model loss depends on more than one resource. A general additive
form is
\begin{equation}
    L(P_1,\ldots,P_m)
    =
    L_{\infty}
    +
    \sum_{i=1}^{m}A_iP_i^{-\alpha_i},
    \label{eq:scaling-general-additive}
\end{equation}
where \(P_i\) is a scaling variable, \(A_i\) is its coefficient, and
\(\alpha_i\) is its exponent. For model size \(N\) and training tokens \(D\),
a conventional formulation is
\begin{equation}
    \boxed{
    L(N,D)
    =
    E+AN^{-\alpha}+BD^{-\beta}
    }
    \label{eq:scaling-additive-nd}
\end{equation}
with positive coefficients \(A,B\), positive exponents \(\alpha,\beta\),
and asymptotic term \(E\).

The limiting cases clarify the interpretation. If model size increases
without bound while data volume is fixed, then
\begin{equation}
    \lim_{N\rightarrow\infty}L(N,D)
    =
    E+BD^{-\beta}.
    \label{eq:scaling-limit-model}
\end{equation}
The model-size contribution vanishes, but the finite-data contribution
remains. Conversely,
\begin{equation}
    \lim_{D\rightarrow\infty}L(N,D)
    =
    E+AN^{-\alpha},
    \label{eq:scaling-limit-data}
\end{equation}
and if both variables increase without bound,
\begin{equation}
    \lim_{\substack{N\rightarrow\infty\\D\rightarrow\infty}}
    L(N,D)=E.
    \label{eq:scaling-joint-limit}
\end{equation}
The additive form is an empirical decomposition of the modeled excess loss;
it does not establish that these are the only possible sources of error.

The additive assumption also differs from a model with explicit interaction
terms. A more general illustrative expression is
\begin{equation}
    L(N,D)
    =
    E+AN^{-\alpha}+BD^{-\beta}
    +GN^{-\mu}D^{-\nu},
    \label{eq:scaling-interaction}
\end{equation}
where the final term represents a possible interaction. The exact functional
form must be supported by the experimental observations and the scientific
hypothesis. The Kaplan-style and Chinchilla-style formulations discussed in
the scaling-law literature encode different assumptions about the interaction
between model size and data volume; their fitted extrapolations can therefore
differ even when both describe a finite set of observations reasonably well
\parencite{Li2025MisFitting,Sengupta2025ScalingSurvey}.

\paragraph{Compute-optimal resource allocation.}
Once a loss function and a computational constraint are specified, the
resource-allocation problem becomes a constrained optimization problem:
\begin{equation}
    \begin{aligned}
        \min_{N,D}\quad & L(N,D)\\
        \text{subject to}\quad & 6ND=C.
    \end{aligned}
    \label{eq:scaling-compute-optimization}
\end{equation}
Substituting \(D=C/(6N)\) into the additive loss in
\cref{eq:scaling-additive-nd} gives
\begin{equation}
    L_C(N)
    =
    E+AN^{-\alpha}
    +B\left(\frac{C}{6}\right)^{-\beta}N^\beta.
    \label{eq:scaling-loss-fixed-compute}
\end{equation}
The second term decreases as \(N\) increases, while the third increases
because a larger model leaves fewer tokens available under the fixed
computational budget. The optimum balances these opposing effects.

Differentiating and setting the derivative to zero gives
\begin{align}
    \frac{dL_C}{dN}
    &=
    -\alpha AN^{-\alpha-1}
    +\beta B\left(\frac{C}{6}\right)^{-\beta}
       N^{\beta-1}=0, \nonumber\\
    \alpha A
    &=
    \beta B\left(\frac{C}{6}\right)^{-\beta}
       N^{\alpha+\beta}, \nonumber\\
    N^{\alpha+\beta}
    &=
    \frac{\alpha A}{\beta B}
       \left(\frac{C}{6}\right)^\beta.
    \label{eq:scaling-optimum-derivation}
\end{align}
Hence the compute-optimal parameter count is
\begin{equation}
    \boxed{
    N^*(C)
    =
    \left(\frac{\alpha A}{\beta B}\right)^{
        \frac{1}{\alpha+\beta}}
    \left(\frac{C}{6}\right)^{
        \frac{\beta}{\alpha+\beta}}
    }
    \label{eq:scaling-optimal-parameters}
\end{equation}
and the corresponding token count, obtained from the compute constraint, is
\begin{equation}
    \boxed{
    D^*(C)
    =
    \left(\frac{\beta B}{\alpha A}\right)^{
        \frac{1}{\alpha+\beta}}
    \left(\frac{C}{6}\right)^{
        \frac{\alpha}{\alpha+\beta}}
    }.
    \label{eq:scaling-optimal-tokens}
\end{equation}
Therefore,
\begin{equation}
    N^*(C)\propto C^{\frac{\beta}{\alpha+\beta}},
    \qquad
    D^*(C)\propto C^{\frac{\alpha}{\alpha+\beta}}.
    \label{eq:scaling-optimal-exponents}
\end{equation}
The optimal allocation depends on the relative loss sensitivities encoded by
\(\alpha\) and \(\beta\), as well as on the coefficients \(A\) and \(B\).

Dividing the two optimal quantities gives the optimal token-to-parameter
ratio:
\begin{equation}
    \boxed{
    \frac{D^*(C)}{N^*(C)}
    =
    \left(\frac{\beta B}{\alpha A}\right)^{
        \frac{1}{\alpha+\beta}}
    \left(\frac{C}{6}\right)^{
        \frac{\alpha-\beta}{\alpha+\beta}}
    }.
    \label{eq:scaling-optimal-ratio}
\end{equation}
If \(\alpha=\beta\), the ratio is independent of compute:
\begin{equation}
    \frac{D^*}{N^*}
    =
    \left(\frac{B}{A}\right)^{\frac{1}{2\alpha}},
    \qquad
    N^*(C)\propto C^{1/2},
    \qquad
    D^*(C)\propto C^{1/2}.
    \label{eq:scaling-constant-ratio}
\end{equation}
If \(\alpha>\beta\), the optimal token-to-parameter ratio increases with
compute; if \(\alpha<\beta\), it decreases. A constant ratio is therefore a
consequence of a particular relationship between fitted exponents, not a
universal property of language-model training.

\paragraph{Kaplan-style and Chinchilla-style conclusions.}
The scaling literature reports different empirical conclusions about how
optimal data volume should grow with model size. The Kaplan-style relationship
is commonly summarized as
\begin{equation}
    D^*\propto (N^*)^{0.74},
    \label{eq:scaling-kaplan-ratio}
\end{equation}
In this relationship, parameter count increases faster than data volume.
A later Chinchilla-style result instead reports approximately proportional
growth:
\begin{equation}
    D^*\propto N^*,
    \qquad
    D^*\approx 20N^*
    \label{eq:scaling-chinchilla-ratio}
\end{equation}
when \(N^*\) is expressed as a parameter count and \(D^*\) as a token count.
The latter is an empirical rule of thumb associated with the studied
conditions, not a universal prescription.

The difference between these conclusions should not be reduced to a claim
that one ratio is correct for every architecture and data regime. The
replication and survey analysis of \textcite{Li2025MisFitting} identifies
training hyperparameters, learning-rate schedules, parameter-counting
conventions, FLOP calculations, checkpoint selection, and fitting procedures
as factors that can alter the estimated relationship. The fitted optimum is
conditional on the experimental design and on the validity of the chosen
functional form.

\begin{figure}[htbp]
    \centering
    % Schematic loss contributions under a fixed-compute constraint.
\begingroup
\setlength{\unitlength}{0.8mm}
\begin{picture}(116,72)
    \put(12,9){\vector(1,0){94}}
    \put(12,9){\vector(0,1){57}}
    \put(108,5){\makebox(0,0)[l]{\(N\)}}
    \put(7,68){\makebox(0,0)[b]{Loss}}
    % Model-capacity contribution: decreasing with N.
    \thinlines
    \qbezier(16,58)(34,22)(99,14)
    \put(77,21){\makebox(0,0)[l]{\(AN^{-\alpha}\)}}
    % Data-limitation contribution: increasing with N at fixed compute.
    \qbezier(16,13)(55,17)(99,57)
    \put(72,53){\makebox(0,0)[l]{\(B(C/6)^{-\beta}N^\beta\)}}
    % Sum: schematic U-shaped objective.
    \thicklines
    \qbezier(16,62)(46,4)(99,60)
    \put(51,8){\makebox(0,0)[t]{\(\displaystyle L_C(N)\)}}
    \put(55,2){\circle*{2.2}}
    \put(55,0){\makebox(0,0)[t]{\(\,N^*\,\)}}
\end{picture}
\endgroup

    \caption{Schematic fixed-compute trade-off. Increasing parameter count
    reduces the model-capacity contribution to loss but increases the
    data-limitation contribution because fewer tokens can be processed.
    Their sum has an interior minimum under the assumptions of the additive
    scaling model.}
    \label{fig:scaling-compute-optimal}
\end{figure}

\paragraph{IsoFLOPs analysis.}
The compute-optimal derivation assumes that the scaling function is known.
In practice, its coefficients and exponents must be estimated. An IsoFLOPs
experiment fixes a computational budget \(C_i\) and evaluates several
configurations satisfying
\begin{equation}
    6ND=C_i,
    \qquad
    D=\frac{C_i}{6N}.
    \label{eq:scaling-isoflops}
\end{equation}
Each configuration is trained and its loss measured. The lowest observed
loss along the fixed-compute profile identifies the empirical optimum for
that budget. Repeating the experiment for several budgets yields observations
of the form
\[
    \{(C_i,N_i^*,D_i^*)\}_{i=1}^{K},
\]
which can be used to estimate how the optimal configuration changes with
compute.

IsoFLOPs analysis directly investigates resource allocation, but it can be
expensive because multiple configurations may need to be trained at each
budget. Intermediate checkpoints and interpolation can reduce the number of
full training runs, although the resulting estimates depend on the
appropriateness of the interpolation procedure. The distinction between
predicting loss for arbitrary configurations and estimating an optimal
resource ratio is important: the former requires a predictive loss surface,
whereas the latter focuses on the minimizing configuration
\parencite{Li2025MisFitting}.

\subsection{Empirical fitting and validation}
\label{subsec:scaling-estimation}

\paragraph{Experimental observations and fitting objective.}
A scaling law is estimated from a dataset of trained and evaluated models.
For the additive form, write
\begin{equation}
    \widehat{L}(N,D;\boldsymbol{\phi})
    =
    E+AN^{-\alpha}+BD^{-\beta},
    \qquad
    \boldsymbol{\phi}=(E,A,B,\alpha,\beta).
    \label{eq:scaling-fitted-function}
\end{equation}
If the \(i\)-th experiment uses \((N_i,D_i)\) and produces measured loss
\(L_i\), the fitting dataset is
\begin{equation}
    \mathcal{S}
    =
    \{(N_i,D_i,L_i)\}_{i=1}^{m}.
    \label{eq:scaling-fitting-dataset}
\end{equation}
The parameter-estimation problem is
\begin{equation}
    \widehat{\boldsymbol{\phi}}
    =
    \arg\min_{\boldsymbol{\phi}}
    \sum_{i=1}^{m}
    \ell\!\left(
        L_i,\widehat{L}(N_i,D_i;\boldsymbol{\phi})
    \right),
    \label{eq:scaling-general-fitting}
\end{equation}
where \(\ell\) is a chosen discrepancy function. The number and distribution
of observations matter. If all experiments use nearly the same
parameter-to-token ratio, it can be difficult to identify the separate
effects of model size and data volume. A narrow range of model sizes also
provides weak evidence for extrapolation to substantially larger models.

Training conditions should be controlled or reported. Architecture, batch
size, learning-rate schedule, optimizer, data composition, training duration,
and evaluation procedure can all affect measured loss. In particular, a
learning rate that is appropriate for a small model may be suboptimal for a
larger model. If optimization quality changes systematically with model
size, the fitted scaling law can mistakenly attribute optimization
differences to scale itself. The survey by \textcite{Li2025MisFitting}
emphasizes the sensitivity of scaling fits to these experimental decisions.

\paragraph{Checkpoint selection and training trajectories.}
A training run provides more than one final observation. Let
\(\boldsymbol{\theta}_i(t)\) denote the parameters of model \(i\) after
processing \(t\) tokens, and let
\begin{equation}
    L_i(t)
    =
    \mathcal{L}_{\mathrm{eval}}(\boldsymbol{\theta}_i(t))
    \label{eq:scaling-learning-trajectory}
\end{equation}
denote its evaluated loss. Intermediate checkpoints can provide observations
at different training durations and can be useful when constructing
fixed-compute profiles. Interpolation between checkpoints can estimate a
loss at an intermediate token count, but this introduces assumptions about
the smoothness of the learning trajectory. The choice between final
checkpoints, multiple checkpoints, smoothed trajectories, and interpolated
observations can affect the fitted law \parencite{Li2025MisFitting}.

\paragraph{Least-squares and robust objectives.}
A common fitting objective is mean squared error:
\begin{equation}
    J_{\mathrm{MSE}}(\boldsymbol{\phi})
    =
    \frac{1}{m}
    \sum_{i=1}^{m}
    \left[
        L_i-\widehat{L}(N_i,D_i;\boldsymbol{\phi})
    \right]^2.
    \label{eq:scaling-mse}
\end{equation}
Squaring residuals prevents positive and negative errors from cancelling and
penalizes large residuals more heavily. Mean absolute error is
\begin{equation}
    J_{\mathrm{MAE}}(\boldsymbol{\phi})
    =
    \frac{1}{m}
    \sum_{i=1}^{m}
    \left|
        L_i-\widehat{L}(N_i,D_i;\boldsymbol{\phi})
    \right|.
    \label{eq:scaling-mae}
\end{equation}
Its linear penalty makes it less sensitive than MSE to individual large
residuals.

The Huber loss combines quadratic and linear behavior. For residual \(r\)
and threshold \(\delta>0\),
\begin{equation}
    \ell_{\delta}(r)
    =
    \begin{cases}
        \frac{1}{2}r^2,
            & |r|\leq\delta,\\[4pt]
        \delta\left(|r|-\frac{\delta}{2}\right),
            & |r|>\delta.
    \end{cases}
    \label{eq:scaling-huber}
\end{equation}
For small residuals it behaves quadratically, while for large residuals it
grows linearly. The threshold \(\delta\) determines the transition. The
scaling-law literature uses variants of MAE, MSE, Huber loss, and custom
objectives; the selected objective should be reported because it changes
the relative influence of observations \parencite{Li2025MisFitting}.

\paragraph{Nonlinear numerical optimization.}
The exponents in a power law are unknown parameters, so fitting is generally
a nonlinear optimization problem. BFGS is a quasi-Newton method that uses
gradient information and an approximation to the Hessian. A Newton-type
update is
\begin{equation}
    \boldsymbol{\phi}_{k+1}
    =
    \boldsymbol{\phi}_{k}
    -
    H_k^{-1}\nabla J(\boldsymbol{\phi}_k),
    \label{eq:scaling-newton-update}
\end{equation}
where \(H_k\) is the Hessian or its approximation. BFGS updates an
approximation to curvature information from successive parameter and
gradient changes. L-BFGS stores only a limited history of these changes,
reducing memory requirements. These methods are used in the scaling-law
literature, including in fitting procedures surveyed by
\textcite{Li2025MisFitting}.

Levenberg--Marquardt is another method for nonlinear least-squares
optimization. If \(\mathbf{r}_k\) is the residual vector and \(J_k\) is its
Jacobian with respect to the fitted parameters, a damped update is
\begin{equation}
    \boldsymbol{\phi}_{k+1}
    =
    \boldsymbol{\phi}_{k}
    -
    \left(J_k^{\mathsf T}J_k+\lambda I\right)^{-1}
    J_k^{\mathsf T}\mathbf{r}_k,
    \label{eq:scaling-levenberg-marquardt}
\end{equation}
where \(\lambda>0\) is a damping parameter and \(I\) is the identity matrix.
Larger damping makes the update more gradient-like, whereas smaller damping
approaches a Gauss--Newton step. Levenberg--Marquardt is used in the
architecture-related scaling analyses of \textcite{Bian2025InferenceEfficient}
and \textcite{Bian2026ArchitectureScaling}.

Initialization is consequential because nonlinear objectives can contain
flat regions, poorly conditioned directions, or multiple local minima.
Grid search, random initializations, and initial values informed by prior
fits are alternative strategies. The replication analysis in
\textcite{Li2025MisFitting} shows that initialization can materially affect
the fitted result. An optimization report should therefore state the
algorithm, objective, initialization, parameter constraints, stopping
criterion, and any model or observation exclusions.

\paragraph{Logarithmic transformations.}
The simple power law \(y=ax^{-\alpha}\) can be linearized as
\begin{equation}
    \log y=\log a-\alpha\log x.
    \label{eq:scaling-log-linear-simple}
\end{equation}
For a loss with an asymptotic term, the corresponding transformation is
\[
    \log(L-E)=\log a-\alpha\log x,
\]
so the asymptotic term must be accounted for. More importantly, minimizing
squared error in the original loss space is not equivalent to minimizing
squared error after a logarithmic transformation:
\begin{equation}
    \sum_i(L_i-\widehat{L}_i)^2
    \neq
    \sum_i(\log L_i-\log\widehat{L}_i)^2.
    \label{eq:scaling-log-objective-difference}
\end{equation}
The transformation changes the error weighting and can give greater
relative influence to errors at smaller loss values. Linearization may be
computationally convenient, but it changes the statistical objective and
should not be assumed to be neutral \parencite{Li2025MisFitting}.

\paragraph{Validation and predictive reliability.}
A function that fits the observations used for estimation is not necessarily
reliable outside that dataset. Residuals are
\begin{equation}
    r_i=L_i-\widehat{L}(N_i,D_i;\widehat{\boldsymbol{\phi}}).
    \label{eq:scaling-residuals}
\end{equation}
Systematic residual patterns can indicate that the assumed functional form
omits an important variable or interaction. The mean squared error and
coefficient of determination are
\begin{align}
    \operatorname{MSE}
    &=
    \frac{1}{m}\sum_{i=1}^{m}r_i^2,
    \label{eq:scaling-validation-mse}\\
    R^2
    &=
    1-
    \frac{\sum_{i=1}^{m}(L_i-\widehat{L}_i)^2}
         {\sum_{i=1}^{m}(L_i-\overline{L})^2},
    \qquad
    \overline{L}=\frac{1}{m}\sum_{i=1}^{m}L_i.
    \label{eq:scaling-r-squared}
\end{align}
The \(R^2\) statistic compares squared prediction error with a baseline that
predicts the sample mean. A high value does not, by itself, establish
accurate extrapolation.

When the objective is to rank candidate architectures, rank agreement can
also matter. For untied ranks \(R_i\) and \(S_i\), Spearman's rank
correlation is
\begin{equation}
    r_s
    =
    1-
    \frac{6\sum_{i=1}^{m}(R_i-S_i)^2}
         {m(m^2-1)}.
    \label{eq:scaling-spearman}
\end{equation}
A value close to one indicates strong agreement between the two rankings.
The inference- and architecture-aware studies use measures including MSE,
\(R^2\), relative prediction error, and Spearman correlation
\parencite{Bian2025InferenceEfficient,Bian2026ArchitectureScaling}.

Confidence intervals, goodness-of-fit measures, and out-of-sample evaluation
provide complementary evidence. A particularly informative extrapolation
test fits a law on smaller models and evaluates its predictions on larger
models not used during fitting. Successful extrapolation over a finite range
is useful evidence, but it does not prove validity at arbitrary scales.

\paragraph{Reproducibility.}
The survey by \textcite{Li2025MisFitting} examines 51 scaling-law studies,
of which 45 quantify scaling trends using power laws. The survey reports
substantial omissions in the description of training, equation variables,
evaluation, fitting procedures, optimization, parameter counting, and FLOP
accounting. It also reports limited availability of analysis code,
checkpoints, and metric scores. These observations show that a compact
equation is not a complete scientific result: the experimental and numerical
pipeline that produced its coefficients must also be documented.

A reproducible scaling study should specify the functional form, definitions
of all variables, parameter-counting convention, FLOP calculation, model
architecture, training data, optimizer, batch size, learning-rate schedule,
checkpoint selection, evaluation metric, fitting objective, optimization
algorithm, initialization strategy, parameter constraints, and validation
procedure. Without these details, independently fitting an equation of the
same nominal form need not reproduce the original conclusion.

\begin{figure}[htbp]
    \centering
    % Schematic IsoFLOPs curves in parameter-count / token-count coordinates.
\begingroup
\setlength{\unitlength}{0.8mm}
\begin{picture}(116,72)
    \put(12,10){\vector(1,0){96}}
    \put(12,10){\vector(0,1){56}}
    \put(110,6){\makebox(0,0)[l]{\(N\)}}
    \put(7,68){\makebox(0,0)[b]{\(D\)}}
    \thinlines
    \qbezier(19,61)(27,36)(83,15)
    \qbezier(19,68)(40,42)(102,20)
    \qbezier(28,68)(56,49)(108,33)
    \put(79,13){\makebox(0,0)[l]{\(C_1\)}}
    \put(98,19){\makebox(0,0)[l]{\(C_2\)}}
    \put(104,33){\makebox(0,0)[l]{\(C_3\)}}
    \thicklines
    \put(45,28){\circle*{2.4}}
    \put(67,35){\circle*{2.4}}
    \put(83,45){\circle*{2.4}}
    \put(45,24){\makebox(0,0)[t]{\( (N_1^*,D_1^*)\)}}
    \put(68,31){\makebox(0,0)[t]{\( (N_2^*,D_2^*)\)}}
    \put(84,41){\makebox(0,0)[t]{\( (N_3^*,D_3^*)\)}}
    \put(14,2){\makebox(90,4)[l]{Each curve fixes approximate training FLOPs.}}
\end{picture}
\endgroup

    \caption{Schematic IsoFLOPs analysis. Each curve represents configurations
    with the same approximate training-compute budget. The minimum-loss
    configuration on each curve defines an empirical compute-optimal point.}
    \label{fig:scaling-isoflops}
\end{figure}

\subsection{Inference-efficient and architecture-aware scaling}
\label{subsec:scaling-inference-architecture}

\paragraph{Training optimality versus deployment optimality.}
The conventional framework minimizes pre-training loss under a training
compute constraint. Deployment introduces additional objectives, including
inference latency, throughput, and memory use. A model that is optimal in
training FLOPs need not be optimal for autoregressive generation.

Parameter count alone does not determine inference latency. In a
decoder-only Transformer, successive layers are executed sequentially during
generation. The number of layers therefore contributes to the sequential
depth of the computation, while hidden dimension affects projection sizes,
matrix multiplication shapes, and memory access. Models with similar
parameter counts can consequently exhibit different inference behavior.
The measurements reported by \textcite{Bian2025InferenceEfficient} show
substantial latency differences between models of comparable size and
illustrate the importance of model shape.

\paragraph{Inference-constrained scaling.}
An inference-aware optimization can be expressed as
\begin{equation}
    \begin{aligned}
        \min_{N,D}\quad & L(N,D)\\
        \text{subject to}\quad
        &N\leq N_C,\\
        &D\leq D_C,\\
        &T_{\mathrm{inf}}\leq T_C,
    \end{aligned}
    \label{eq:scaling-inference-constrained}
\end{equation}
where \(N_C\) and \(D_C\) are the available model-size and training-data
budgets, and \(T_C\) is the maximum permitted inference latency. The
latency \(T_{\mathrm{inf}}\) must be measured or modeled under a specified
hardware, batch-size, sequence-length, and decoding configuration; it is
not a universal property of parameter count alone.

\textcite{Bian2025InferenceEfficient} introduce the aspect ratio
\begin{equation}
    R=\frac{d_{\mathrm{model}}}{n_{\mathrm{layers}}},
    \label{eq:scaling-aspect-ratio}
\end{equation}
where \(d_{\mathrm{model}}\) is the hidden dimension and
\(n_{\mathrm{layers}}\) is the number of Transformer layers. Their
inference-efficient loss formulation is
\begin{equation}
    \boxed{
    L(N,D,R)
    =
    \left(E+AN^{-\alpha}+BD^{-\beta}\right)
    \left(1+\epsilon R^\gamma\right)
    }
    \label{eq:scaling-inference-efficient-law}
\end{equation}
where \(\epsilon\) and \(\gamma\), together with the conventional scaling
coefficients, are estimated from observations. The first factor represents
the conventional size-and-data relationship, while the second introduces
an aspect-ratio-dependent adjustment.

The authors train 63 models across different parameter counts, data budgets,
and architectural configurations. Their experiments report that models with
similar parameter counts can have different inference latencies and that
wider, shallower configurations can offer throughput advantages under the
studied settings. They use the fitted relationship together with latency
measurements to select candidate architectures. Their reported Morph-1B
result demonstrates an inference-latency improvement relative to comparable
open-source models under the paper's evaluation conditions. These findings
support the broader principle that model selection should account for
deployment constraints rather than optimizing pre-training loss in isolation
\parencite{Bian2025InferenceEfficient}.

\paragraph{Conditional architecture-aware scaling.}
Architecture is not merely a fixed background condition. Hidden dimension,
layer count, and the allocation of parameters between attention and
feed-forward networks are also optimization variables. In a decoder-only
Transformer, let
\begin{equation}
    r_{\mathrm{MLP/attn}}
    =
    \frac{N_{\mathrm{MLP}}}{N_{\mathrm{attn}}}
    \label{eq:scaling-mlp-attention-ratio}
\end{equation}
denote the ratio of feed-forward parameters to attention parameters.
Grouped-query attention (GQA) changes the relationship between query heads
and key-value heads and can therefore affect key-value-cache memory and
inference throughput.

The architecture-aware study of \textcite{Bian2026ArchitectureScaling}
reports approximately U-shaped relationships between predictive loss and
normalized hidden dimension, and between loss and the MLP-to-attention
ratio. This behavior differs from a monotonically decreasing power law:
under a fixed parameter budget, an architectural variable can have an
interior optimum. The paper uses a representative calibration form
\begin{equation}
    f(x)=c_0+c_1\log x+\frac{c_2}{x},
    \label{eq:scaling-architecture-u-shape}
\end{equation}
where \(x>0\). Its derivative is
\[
    f'(x)=\frac{c_1}{x}-\frac{c_2}{x^2}.
\]
For positive \(c_1,c_2\), the stationary point is
\begin{equation}
    x^*=\frac{c_2}{c_1},
    \qquad
    f''(x^*)=\frac{c_1^3}{c_2^2}>0,
    \label{eq:scaling-architecture-minimum}
\end{equation}
so it is a local minimum.

The normalized hidden dimension is defined in the study as
\begin{equation}
    x_d=
    \frac{d_{\mathrm{model}}}
         {\sqrt{N_{\mathrm{nonembed}}}},
    \label{eq:scaling-normalized-hidden-dimension}
\end{equation}
where \(N_{\mathrm{nonembed}}\) is the non-embedding parameter count. Let
\(r=r_{\mathrm{MLP/attn}}\). The conditional calibration is represented as
\begin{align}
    L(x_d,r\mid N,D)
    ={}&
    \left(a_0+a_1\log x_d+\frac{a_2}{x_d}\right)
    \nonumber\\
    &\times
    \left(b_0+b_1\log r+\frac{b_2}{r}\right)
    L_{\mathrm{opt}}(N,D),
    \label{eq:scaling-conditional-architecture-law}
\end{align}
where
\begin{equation}
    L_{\mathrm{opt}}(N,D)
    =
    E+AN^{-\alpha}+BD^{-\beta}.
    \label{eq:scaling-architecture-reference-loss}
\end{equation}
The coefficients are fitted to experimental observations. The multiplicative
form assumes that the two architectural calibration factors can be separated;
it is a modeling choice rather than a general property of Transformer
architectures.

Architecture selection can then be written as an efficiency optimization
subject to a loss constraint:
\begin{equation}
    \begin{aligned}
        \max_{\mathcal{P}}\quad & I_N(\mathcal{P})\\
        \text{subject to}\quad
        &L(\mathcal{P}\mid N,D)\leq L_t,
    \end{aligned}
    \label{eq:scaling-architecture-optimization}
\end{equation}
where \(\mathcal{P}\) denotes an architectural configuration,
\(I_N(\mathcal{P})\) is its inference-efficiency measure, and \(L_t\) is
the maximum acceptable loss. GQA can be treated as a discrete architectural
choice and searched over feasible configurations rather than forced into
the same continuous calibration function.

The study trains more than 200 models, spanning approximately 80 million
to 3 billion non-embedding parameters and training-data budgets from
8 billion to 100 billion tokens. It reports up to 42\% greater inference
throughput and 2.1\% higher average downstream accuracy than its specified
LLaMA-3.2 baselines under the reported conditions. The authors also show
that using models closer to the target scale can improve prediction accuracy.
These results reinforce that the scale range used for fitting remains
important even when architecture is explicitly modeled
\parencite{Bian2026ArchitectureScaling}.

\begin{figure}[htbp]
    \centering
    % Schematic contrast between deep/narrow and shallow/wide configurations.
\begingroup
\setlength{\unitlength}{0.8mm}
\begin{picture}(150,57)
    \put(3,45){\makebox(61,5)[c]{\textbf{Deep / narrow}}}
    \put(82,45){\makebox(64,5)[c]{\textbf{Shallow / wide}}}
    % Deep/narrow stack.
    \put(8,32){\framebox(18,8){Layer 1}}
    \put(8,21){\framebox(18,8){Layer 2}}
    \put(8,10){\framebox(18,8){Layer 3}}
    \put(17,32){\vector(0,-1){3}}
    \put(17,21){\vector(0,-1){3}}
    \put(30,24){\makebox(33,7)[l]{More sequential stages}}
    % Shallow/wide stack.
    \put(88,27){\framebox(43,13){Wide layer 1}}
    \put(88,10){\framebox(43,13){Wide layer 2}}
    \put(109,27){\vector(0,-1){4}}
    \put(135,24){\makebox(12,7)[l]{Fewer stages}}
    \put(3,1){\makebox(143,5)[l]{At a fixed parameter budget, shape can change latency and predictive loss.}}
\end{picture}
\endgroup

    \caption{Schematic architecture-aware scaling. At a fixed parameter
    budget, different allocations between depth and width can have different
    inference costs and predictive losses; an architecture-aware objective
    searches for efficient configurations subject to a loss constraint.}
    \label{fig:scaling-architecture}
\end{figure}

\subsection{Joint scaling for mixture-of-experts models}
\label{subsec:scaling-moe}

\paragraph{Stored parameters and active parameters.}
In a dense Transformer, the same feed-forward parameters are applied to each
token. A mixture-of-experts model replaces a dense feed-forward component
with multiple expert networks and a routing mechanism. The router selects
which expert or experts process a token. In the top-one routing configuration
considered in the Switch-style setting, a token is routed to one expert.

Let \(N_{\mathrm{total}}\) be the total number of stored parameters and
\(N_{\mathrm{act}}\) the number of parameters activated for a token. For a
dense model, these quantities are approximately equal under a consistent
counting convention. For an MoE model,
\begin{equation}
    N_{\mathrm{act}}<N_{\mathrm{total}},
    \label{eq:scaling-moe-active-total}
\end{equation}
because only a subset of expert parameters participates in a token's
computation. Active parameters are relevant to arithmetic work per token,
whereas total parameters are important for weight storage. The number of
experts, denoted by \(E\), is an additional architectural variable.

\paragraph{Joint MoE scaling law.}
\textcite{Ludziejewski2025JointMoE} extend the Chinchilla-style relationship
to include active parameters, training tokens, and the number of experts.
Their joint formulation is
\begin{align}
    L(N_{\mathrm{act}},D,\widehat{E})
    ={}&
    a\widehat{E}^{\delta}
    N_{\mathrm{act}}^{\alpha+\gamma\ln(\widehat{E})}
    \nonumber\\
    &+
    b\widehat{E}^{\omega}
    D^{\beta+\zeta\ln(\widehat{E})}
    +c,
    \label{eq:scaling-joint-moe}
\end{align}
where \(\widehat{E}\) is the transformed expert-count variable defined by
the study, and \(a,b,c,\delta,\alpha,\beta,\gamma,\omega,\zeta\) are fitted
parameters. The effective exponents are
\begin{equation}
    \mu(E)=\alpha+\gamma\ln(\widehat{E}),
    \qquad
    \nu(E)=\beta+\zeta\ln(\widehat{E}).
    \label{eq:scaling-moe-effective-exponents}
\end{equation}
Thus, the influence of active parameters and data volume can change with
the expert count. For a fixed expert configuration, the expression retains
a Chinchilla-style dependence on active parameters and training tokens.

The transformation \(\widehat{E}\) is part of the authors' specified
functional form and should be implemented according to their definition;
it should not be replaced by an assumed transformation when reproducing
their fit. The formulation is empirical and its coefficients and exponents
must be estimated from observations.

\paragraph{Compute- and memory-constrained MoE scaling.}
Under the conventional approximation used for the studied sparse models,
training computation is expressed using active parameters:
\begin{equation}
    F_{\mathrm{train}}
    \approx 6N_{\mathrm{act}}D.
    \label{eq:scaling-moe-flops}
\end{equation}
The compute-optimal problem is therefore
\begin{equation}
    \begin{aligned}
        \min_{N_{\mathrm{act}},D,E}\quad
        &L(N_{\mathrm{act}},D,E)\\
        \text{subject to}\quad
        &6N_{\mathrm{act}}D=F,
    \end{aligned}
    \label{eq:scaling-moe-compute-opt}
\end{equation}
where \(F\) is the available training-compute budget. Unlike the dense
formulation, the expert count now influences the optimal balance between
active parameters and tokens.

The results reported by \textcite{Ludziejewski2025JointMoE} show that, for
the studied settings and fixed training-compute budget, increasing the
number of experts leads to an optimum with fewer active parameters and more
training tokens. Their reported configurations for a budget of
\(10^{20}\) FLOPs are shown in \cref{tab:scaling-moe-configurations}.

\begin{table}[htbp]
    \centering
    \caption{Reported compute-optimal MoE configurations at a training
    budget of \(10^{20}\) FLOPs. Values are reproduced from the results
    reported by \textcite{Ludziejewski2025JointMoE}.}
    \label{tab:scaling-moe-configurations}
    \begin{tabularx}{\textwidth}{c>{\centering\arraybackslash}X>{\centering\arraybackslash}X}
        \toprule
        Number of experts & \shortstack{Optimal active\\parameters} & \shortstack{Optimal training\\tokens}\\
        \midrule
        1  & 1.7B  & 9.7B\\
        2  & 1.5B  & 11.4B\\
        4  & 1.2B  & 13.9B\\
        8  & 990M  & 17B\\
        16 & 810M  & 20.7B\\
        \bottomrule
    \end{tabularx}
\end{table}

The optimal ratio of tokens to active parameters changes with expert
count. A single fixed ratio cannot describe every MoE configuration.

Training compute is not the only constraint. The total parameter count
determines weight-storage requirements, and autoregressive inference also
requires memory for the key-value cache. For a decoder with \(N_{\mathrm{blocks}}\)
blocks, hidden dimension \(d_{\mathrm{model}}\), and \(T\) cached tokens, a
simplified cache-element count is proportional to
\begin{equation}
    M_{\mathrm{KV}}
    \propto
    2T N_{\mathrm{blocks}}d_{\mathrm{model}},
    \label{eq:scaling-moe-kv-cache}
\end{equation}
where the factor of two accounts for keys and values. The exact memory in
bytes additionally depends on batch size, the number of key-value heads,
storage precision, and implementation details. Under a simplified
parameter-equivalent accounting, the combined constraint can be represented
as
\begin{equation}
    N_{\mathrm{total}}
    +
    2T N_{\mathrm{blocks}}d_{\mathrm{model}}
    \leq M,
    \label{eq:scaling-moe-memory}
\end{equation}
where \(M\) is the corresponding memory budget. This expression is a
simplified representation of the weight and cache requirements, not a
hardware-independent byte-accurate memory model.

The MoE analysis therefore distinguishes three quantities that must not be
collapsed into one: total stored parameters, active parameters per token,
and memory required by the key-value cache. The study argues that MoE models
can be memory-efficient under the investigated constraints even though
their total parameter count is larger than that of a corresponding dense
model. This conclusion depends on the studied expert configurations and
accounting assumptions; FLOPs alone do not capture all communication,
memory-access, or hardware execution costs
\parencite{Ludziejewski2025JointMoE}.

\begin{figure}[htbp]
    \centering
    % Schematic distinction between dense and mixture-of-experts parameter use.
\begingroup
\setlength{\unitlength}{0.8mm}
\begin{picture}(150,60)
    \put(5,51){\makebox(58,5)[c]{\textbf{Dense feed-forward block}}}
    \put(85,51){\makebox(60,5)[c]{\textbf{Mixture of experts}}}
    \put(8,34){\framebox(45,11){Shared parameters}}
    \put(30,25){\vector(0,1){8}}
    \put(30,17){\makebox(0,0)[c]{\(\text{each token}\)}}
    \put(88,35){\framebox(18,9){Expert 1}}
    \put(111,35){\framebox(18,9){Expert 2}}
    \put(134,35){\framebox(12,9){\(\cdots\)}}
    \put(105,23){\framebox(25,8){Router}}
    \put(117,31){\vector(0,1){3}}
    \put(117,14){\vector(0,1){8}}
    \put(117,8){\makebox(0,0)[c]{Token}}
    \put(5,2){\makebox(140,5)[l]{MoE stores many expert parameters but activates a selected subset per token.}}
\end{picture}
\endgroup

    \caption{Schematic distinction between dense and mixture-of-experts
    parameter accounting. A dense model activates its shared parameters for
    each token, whereas an MoE model stores multiple experts and activates
    only a selected subset for a token.}
    \label{fig:scaling-moe}
\end{figure}

\subsection{Limitations, synthesis, and transition}
\label{subsec:scaling-limitations-transition}

The scaling relationships developed in this section provide a quantitative
framework for reasoning about model size, training data, computational
expenditure, inference efficiency, and architectural allocation. Their
predictive validity is nevertheless conditional. A scaling law fitted for
one architecture should not automatically be transferred to another
architecture, because equal parameter counts can correspond to different
hidden dimensions, layer counts, attention configurations, and allocations
between attention and feed-forward components. The same caution applies to
data: equal token counts do not imply equal information content, domain
coverage, noise, or repetition.

Extrapolation is another fundamental limitation. A power law fitted over a
restricted range can fail when the system enters a different computational
or statistical regime. Piecewise or broken-scaling formulations have been
proposed to represent changes in scaling behavior, but such extensions
require evidence for a transition and additional parameters to estimate.
Likewise, predictive loss does not uniquely determine downstream accuracy.
The mapping from pre-training loss to task performance can differ across
tasks, and downstream evaluation remains necessary when task performance
is the actual objective \parencite{Sengupta2025ScalingSurvey,
Bian2025InferenceEfficient}.

The conventional dense-model formulation begins with
\(L(N,D)=E+AN^{-\alpha}+BD^{-\beta}\) and the approximate computational
constraint \(C\approx6ND\). Their combination yields compute-optimal
parameter and token allocations. Empirical estimation then requires a
carefully designed collection of training runs, explicit definitions of
variables, numerical fitting, and validation beyond the observations used
for optimization. Inference-aware and architecture-aware formulations add
latency and architectural variables, while MoE scaling distinguishes stored
parameters from active parameters and introduces memory constraints.

The unifying principle is that scale is multidimensional. Increasing one
resource changes the value and cost of the others, and the best configuration
depends on the objective being optimized. Scaling laws are useful precisely
because they turn these interactions into testable and predictive
relationships, but their conclusions remain tied to their assumptions,
measurements, and experimental domains. This foundation prepares the
subsequent discussion of how training objectives, optimization procedures,
and architectural choices interact in the construction of increasingly
capable language models.
